% =========================================================================
% "Minimal dual-oscillator phase model" Methods subsection (reduced model model).
% Written to be pasted into the manuscript's Materials and Methods; adapt
% equation numbering and citation keys to house style. Parameter tables,
% the full fitting objective, numerical checks and the sensitivity analysis
% are in supplementary_model.tex.
% Source code: code/phase_model.py.
% =========================================================================
\section*{Methods: minimal dual-oscillator phase model}

We modeled the morning (M) and evening (E) clock-neuron transcriptional
rhythms as two coupled phase oscillators, $\phi_M(t)$ and $\phi_E(t)$,
with a \textit{tim} transcriptional peak defined at $\phi = \pi$. Time
$t$ is in hours, with zeitgeber time (ZT) 0 at lights-on. The phases
evolved as
\begin{equation}
\begin{aligned}
\dot\phi_M &= \omega + K\,L(t)\sin\bigl(\theta(t)-\phi_M\bigr)
              + \kappa\sin(\phi_E-\phi_M) + g_M\,\gamma\,S(t), \\
\dot\phi_E &= \omega + K\,L(t)\sin\bigl(\theta(t)-\phi_E\bigr)
              - \kappa\sin(\phi_E-\phi_M) - g_E\,\gamma\,S(t).
\end{aligned}
\end{equation}
The four terms are, in order, the intrinsic rhythm
($\omega = 2\pi/24\ \mathrm{h}^{-1}$), entrainment by light, mutual
coupling, and a PDF-dependent drive. The two oscillators differ only in
the sign of the PDF-dependent term, so that without PDF signaling their
phases are identical and their peaks coincide.

Light was a square wave, $L(t) = 1$ for $(t \bmod 24) < P$ and $0$
otherwise, with $P$ the photoperiod in hours. In the light, each
oscillator is attracted with strength $K$ toward a reference phase
$\theta(t)$ that reaches the peak phase $\pi$ at ZT $B(P)$:
\begin{equation}
\begin{aligned}
\theta(t) &= \omega t + \pi - \omega\,B(P), \qquad
B(P) = \max(8,\,P_c+1) + (\lambda-1)\,F(P), \\
P_c &= \min(P,\,15), \qquad x = \max(0,\,P_c-9), \qquad
F(P) = \frac{x^2}{x^2+1}.
\end{aligned}
\end{equation}
The reference peak is therefore at ZT~8 on days of 7~h or shorter and
1~h after lights-off at 7--9~h; on longer days its delay after lights-off
grows toward $\lambda$ hours as the activation function $F(P)$ rises from
0 toward 1, and it no longer changes beyond 15~h of light. These
landmarks were fixed from the measured peak times.

The effective phase-separating action of PDF was
\begin{equation}
S(t) = L(t)\,\max\bigl(0,\cos\theta(t)\bigr)\,F(P),
\end{equation}
which is present only in the light, during the subjective morning of the
reference cycle, and at photoperiods longer than 9~h. It advances M and
delays E with equal strength $\gamma$, consistent with the opposite phase
shifts produced by cell-specific loss of the PDF receptor (PDFR). The
flags $(g_M, g_E)$ switch the drive on or off in each cell group: $(1,1)$
for wild type, $(0,0)$ for loss of PDF, $(0,1)$ for M-specific and
$(1,0)$ for E-specific loss of PDFR. The coupling $\kappa$ pulls the two
phases together in every genotype; in constant darkness (DD), where
$L = S = 0$, it is the only term acting on the M--E gap, which then
decays with a time constant of $1/(2\kappa)$.

The equations were integrated numerically
(\texttt{scipy.integrate.odeint}), with lights-on and lights-off supplied
as integration breakpoints, and peak times were taken as the times at
which each phase crossed $\pi$ once the rhythm had settled (60 simulated
days). Four parameters ($K$, $\lambda$, $\gamma$, $\kappa$) were
estimated by bounded nonlinear least squares
(\texttt{scipy.optimize.least\_squares}). The targets were the bootstrap
mean M and E peak times at 11 wild-type photoperiods (10,000 paired,
stratified resamples of the 20 brains at each ZT), weighted by their
bootstrap covariance; the M and E peak times of the three mutant
genotypes at 12L:12D; and the M--E gap on days 1 and 4 after transfer
from 12L:12D to DD. The resulting fit ($K=0.073$, $\gamma=0.108$ and
$\kappa=0.0062~\mathrm{rad\,h^{-1}}$; $\lambda=3.01$~h) reproduced the
wild-type peak times with a root-mean-square error of 0.23~h and the
mutant peak times to within 0.35~h. The complete fitting objective,
parameter and target tables, numerical checks and sensitivity analysis
are given in the Supplementary Information. Model code, fitted
parameters, and the experimental targets used for fitting are available
at [repository URL / data availability statement].
